% !TEX TS-program = xelatex
% !TEX encoding = UTF-8 Unicode

\documentclass{resume_zh}
\renewcommand{\baselinestretch}{1.0}
\usepackage{zh_CN-Adobefonts_external}
\usepackage{linespacing_fix}
\usepackage{enumitem}

\begin{document}
\pagenumbering{gobble}

\name{钱煦}

\basicInfo{
  \email{drqianxu@163.com} \textperiodcentered\
  \phone{(+86) 156-0061-6741} \textperiodcentered\
  \github[DrXuQian]{https://github.com/DrXuQian}}

\section{个人简介}
8年以上 AI 推理优化经验（NPU/GPU），覆盖模型量化、CUDA Kernel 开发、编译器设计及软硬件协同设计全栈。

\section{教育背景}
\datedsubsection{\textbf{中国科学院大学}，电子工程博士}{2013 -- 2018}
优秀毕业生
\datedsubsection{\textbf{东南大学}，电子工程学士}{2009 -- 2013}

\section{专业技能}
\begin{itemize}[parsep=0.5ex]
  \item \textbf{编程语言}: Python, C++, CUDA
  \item \textbf{框架与工具}: PyTorch, TensorRT-LLM, vLLM, CUTLASS, MLIR, Nsight Systems/Compute
  \item \textbf{专业领域}: 模型量化、Kernel 优化、NPU/GPU 架构设计、编译器开发
\end{itemize}

\section{工作经历}

\datedsubsection{\textbf{阿里巴巴 PTG}，端侧计算组 AI 软件架构师}{2024 Q4 -- 至今}
\textbf{VLM/VLA/MLLM 端侧推理优化}

面向端侧设备的 VLM/VLA/MLLM 推理加速与模型优化，部署于自研 GPGPU AI 加速器。

\begin{itemize}

\item \textbf{模型量化与压缩}
\begin{itemize}
  \item \textbf{Qwen-VL}: 采用 AWQ、SmoothQuant、GPTQ、SpinQuant 量化至 FP4（权重+激活），精度优于 W4A16 baseline，同时 Prefill 延迟降低 50\%。集成 Eagle speculative decoding，Decode 阶段加速 $\sim$4.7$\times$。
  \item \textbf{Orion}（VLA 模型）: 组合 Spin + Scale + SVD 分解 + 混合精度，保持 FP4 模型精度，延迟较 FP8 baseline 降低 $>$15\%。
  \item \textbf{RDQuant} -- 率失真最优混合精度量化：提出基于率失真优化的逐输出通道混合精度框架，通过 Lagrangian 松弛在给定比特预算下最优分配比特宽度（NVFP4/FP8/FP16）。设计混合精度 GEMM 执行 pipeline，包括precision-aware channel regrouping、weight repacking 及 multi-format kernel dispatch。在 Qwen3-4B 上以 $\sim$5.3 bits/weight 超越 BF16 baseline perplexity。已集成至 vLLM，在 NVIDIA RTX 5090 上实现 1.9$\times$ 加速。
  \item \textbf{面向硬件 Compression IP 的 Entropy-Aware 压缩}：开发面向后端硬件 compression IP 的权重与 KV cache 压缩方案。权重侧提出逐 block FP8 scale 优化，联合最小化 MSE 与 codeword entropy，NVFP4 精度无损下实现 $\sim$12\% 压缩。KV cache 侧提出基于正交旋转和 Lloyd-Max 量化的低熵 FP16 表示，实现等效 $\sim$3-bit 的硬件透明压缩，无需定制 dequantization kernel。
  \item 对 Microscaling FP4 和 NVFP4 进行软硬件联合评估，指导量化硬件路线图。
\end{itemize}

\item \textbf{Kernel 开发与算子优化}
\begin{itemize}
  \item \textit{自动驾驶算子}: 优化 DeformableAttention、GridSample、Voxelization、BevPoolv2、CoordTransformMatrix 等性能关键算子。DeformableAttention 延迟降至 MMCV baseline 的 $\sim$1/6。
  \item \textit{大模型算子}:
  \begin{itemize}
    \item \textbf{W4A16 推理优化}: 针对 Decode 阶段特定 shape（尤其 $M\!=\!1$）优化 W4A16 GEMV kernel，bandwidth utilization 达 $\sim$80\%。在 Marlin 后端实现 Fused SwiGLU-FFN kernel，重排 gate/up 权重实现 register 级 GeLU + Mul 融合。
    \item 将 \textbf{MoE W4A16} kernel 集成至内部 TensorRT-LLM pipeline，包括offline weight repacking 与 graph integration，支持 Qwen3-MoE 端到端部署。
    \item \textbf{FP8 GEMM} 优化: 在 CUTLASS 中添加 M 维度对齐至 48 整数倍的支持，启用 PTX 层 1.5$\times$ throughput 特性。修复 TMA store partial-write 导致的 2$\times$ write traffic 问题。
  \end{itemize}
\end{itemize}

\item \textbf{跨平台 Profiling、性能分析与硬件协同设计}
\begin{itemize}
  \item 使用 Nsight Systems 和 Nsight Compute（内部版本）对 Qwen2.5-VL、Qwen3-VL、Qwen3-Omni 在 NVIDIA Thor 和自研 GPGPU 上进行端到端延迟 profiling 与关键算子 bottleneck 分析。
  \item 针对 VAD-Base 中 Deformable Attention 进行bottleneck 分析，导出 \textbf{double ALU} 硬件需求，定量评估性能收益并交付 GPGPU 团队。
  \item 分析大模型 Prefill 阶段 Flash Attention 硬件优化选项，定量评估 \textbf{double SFU} 收益并作为 SoC 设计输入。
\end{itemize}

\end{itemize}

\datedsubsection{\textbf{Intel}，Movidius NPU 团队 AI 软件架构师}{2022 Q1 -- 2024 Q4}
\textbf{NPU 软件架构设计}

参与多代 Intel Movidius NPU 的软硬件架构协同设计，聚焦原型编译器开发、高级编译技术及软硬件协同优化。

\begin{itemize}

\item \textbf{原型编译器开发}: 主导开发原型编译器，对接硬件仿真器实现精准的模型性能模拟。
\begin{itemize}
  \item 评估未来 NPU 硬件变更对软件栈的影响，识别 PPA 最优的架构演进路径。
  \item 定义关键性能指标，分析客户负载（Microsoft、HP 等）性能瓶颈，指导编译器特性规划。
  \item 支持 800+ 模型，Intel Open Model Zoo 覆盖率 90\%+。
\end{itemize}

\item \textbf{高级编译技术}: 与产品编译器和 DirectML 团队协作，探索并验证前沿编译优化。
\begin{itemize}
  \item 通过混合精度、FlashAttention、task pipelining和内存高效 FFN tiling 优化 LLM 推理。
  \item 开发自动 layer fusion 算法和 double buffering 调度。
  \item 设计计算/内存资源分配和 prefetch 调度算法，实现计算与通信重叠。
  \item 在基于 MLIR 的产品编译器上完成 POC 验证。
\end{itemize}

\item \textbf{NPU 软硬件协同设计}: 参与硬件定义和架构探索，确保 PPA 最优的软硬件对齐。
\begin{itemize}
  \item 评估 MAC 阵列大小、SRAM 带宽、PPE 宽度、累加器上下文大小和 DSP 指令宽度的设计权衡。
  \item 提出 Transformer 动态 shape 硬件支持方案，支持高效 MoE 执行。
  \item 定义 tile 互连策略（unicast/multicast/broadcast），平衡延迟、带宽和可扩展性。
  \item 参与存储子系统设计，包括 scratchpad SRAM、cache 层级和 DMA 策略。
\end{itemize}

\end{itemize}

\datedsubsection{\textbf{Intel}，深度学习软件工程师}{2018 Q4 -- 2022 Q1}
\textbf{NPU 模型压缩}

主导 NPU 模型压缩算法开发，覆盖量化、剪枝、AutoML 和神经架构搜索。

\begin{itemize}
\item \textbf{低比特量化}: 开发混合精度量化流水线（PACT、SAWB、HAQ、HAWQ）及非线性量化方法。
\item \textbf{MLPerf 模型剪枝}: 交付 MLPerf 推理基准的超压缩模型（Resnet-50、Mobilenet、SSD-Resnet34、Bert）。Resnet-50 在 1\% 精度损失内实现 84\% weight sparsity和 65\% activation sparsity，speedup $>$2$\times$。
\item \textbf{AutoML 量化}: 基于强化学习开发通道级 Hessian 感知量化方法，top-1 精度超越同期最优方法 $\sim$1\%。
\item \textbf{神经架构搜索}: 面向 NPU 设计优化网络 --- 超分辨率模型比 EDSR 快 6$\times$，PSNR 提升 7\%；分类模型比 Mobilenet-V2 快 10\%，top-1 精度提升 5\%。
\end{itemize}

\datedsubsection{\textbf{Intel}，机器学习工程师}{2018 Q3 -- 2018 Q4}
\textbf{基于机器学习的长城虚拟修复}

与中国文物保护基金会合作，应用 3D-GAN 和纹理合成（GAN）进行长城结构 AI 修复。

\datedsubsection{\textbf{中国科学院}，博士研究生}{2013 Q3 -- 2018 Q2}
\begin{itemize}
  \item \textbf{自研 WiFi 芯片软件设计}: 参与两款 802.11 芯片流片，设计 MAC 层并编写 Linux 驱动。
  \item \textbf{基于 ML 的 SNR 预测}: 开发 MAC 层 SNR 预测的 ML 方法，吞吐量提升 $>$20\%，并在 LSTM 上应用模型压缩。
\end{itemize}

\section{论文发表}
\begin{enumerate}
  \item \textbf{Xu Qian}, Victor Li, Crews Darren S, ``Neural Architecture Search for Intel Movidius VPU'', \textit{arxiv 2305.03739}
  \item \textbf{Xu Qian}, Victor Li, Crews Darren S, ``Auto Hessian Aware Channel-wise Quantization of Neural Networks'', \textit{EMC2 2020}
  \item Charles Qi, Yi Wang, et al., \textbf{Xu Qian}, ``VPU-EM: An Event-based Modeling Framework to Evaluate NPU Performance and Power Efficiency at Scale'', \textit{arxiv 2303.10271}
  \item \textbf{Xu Qian}, Bin Wu, Tianchun Ye, ``QoS-aware A-MPDU retransmission scheme for 802.11 n/ac/ad WLANS'', \textit{IEEE Communications Letters 21 (10), 2290-2293}
  \item \textbf{Xu Qian}, Bin Wu, Tianchun Ye, ``Enhanced aggregation scheduler design in industrial 802.11 devices'', \textit{Electronics Letters 53 (15), 1073-1075}
\end{enumerate}

\section{专利}
\begin{enumerate}
  \item Crews Darren S, et al., \textbf{Xu Qian}, ``Methods and apparatus to accelerate convolution'', \textbf{WO2023044707A1}
  \item \textbf{Xu Qian}, et al., ``Apparatus and method for RL based post-training sparsification'', \textbf{WO2023082278A1}
  \item \textbf{Xu Qian}, et al., ``Apparatus for accelerating computation of process engine'', \textbf{WO2023092383A1}
  \item Peiqing Jiang, et al., \textbf{Xu Qian}, ``Executing MatMul By Performing Conv With DNN Accelerator'', \textbf{WO2025184850A1}
  \item Haiyun Hong, et al., \textbf{Xu Qian}, ``Reshaping Conv Based On Config Of DNN Accelerator'', \textbf{WO2025189339A1}
  \item Haiyun Hong, et al., \textbf{Xu Qian}, ``Streamline Processing Graph for NN Execution'', \textbf{WO2026000674A1}
  \item Yuanyuan Li, \textbf{Xu Qian}, et al., ``Compact Config Descriptor for NN HW Accelerator'', \textbf{WO2026000671A1}
  \item Dinakar Kondru, \textbf{Xu Qian}, et al., ``Flexible Inter-Tile Comm in Multi-Tile NN Accelerator'', \textbf{WO2026000675A1}
  \item[+] 另有 3 项专利已提交，审核中。
\end{enumerate}

\end{document}
